\chapter{TỔNG QUAN}
\section{Lý do chọn đề tài}
Giáo dục đại học là một khoản đầu tư dài hạn, trong khi lựa chọn không phù hợp có thể dẫn đến giảm động lực học tập, chuyển ngành, kéo dài thời gian tốt nghiệp hoặc làm việc trái chuyên môn. Ở cấp độ xã hội, sự lệch pha giữa sở thích, năng lực và nhu cầu nhân lực làm giảm hiệu quả phân bổ nguồn lực. Đối với cơ sở đào tạo, hiểu đúng tiêu chí lựa chọn giúp hoạt động tuyển sinh chuyển từ truyền thông đại trà sang cung cấp thông tin có ích và có trách nhiệm.

Các nghiên cứu quốc tế cho thấy lựa chọn trường chịu ảnh hưởng đồng thời bởi đặc điểm học sinh, người có ảnh hưởng và đặc tính của cơ sở đào tạo \citep{chapman1981,hossler1987}. Lý thuyết hành vi có kế hoạch nhấn mạnh thái độ, chuẩn chủ quan và khả năng kiểm soát cảm nhận trong việc hình thành ý định \citep{ajzen1991}. Bối cảnh Việt Nam bổ sung vai trò của cơ hội việc làm, danh tiếng, học phí, gia đình và nguồn thông tin \citep{le2022,hai2023}. Tuy nhiên, bằng chứng riêng cho ĐBSCL và cách kết nối quyết định chọn trường với chọn ngành vẫn cần được củng cố.

Theo cổng tuyển sinh chính thức, ĐHCT cung cấp nhiều lựa chọn ngành và chương trình đào tạo trong mùa tuyển sinh 2025 \citep{ctu2025}. Sự đa dạng tạo thêm cơ hội nhưng đồng thời làm tăng yêu cầu tư vấn phù hợp. Đề tài vì thế có ý nghĩa cả về học thuật lẫn quản trị tuyển sinh tại ĐHCT và các cơ sở giáo dục trong vùng.

\section{Mục tiêu nghiên cứu}
\subsection{Mục tiêu chung}
Xác định, đo lường và giải thích các yếu tố ảnh hưởng đến quyết định chọn trường/ngành học của học sinh THPT tại các tỉnh ĐBSCL; từ đó đề xuất hàm ý giúp nâng cao chất lượng hướng nghiệp và truyền thông tuyển sinh.

\subsection{Mục tiêu cụ thể}
\begin{enumerate}[label=\arabic*)]
\item Hệ thống hóa cơ sở lý thuyết và xây dựng mô hình nghiên cứu phù hợp với học sinh THPT tại ĐBSCL.
\item Kiểm định độ tin cậy, giá trị thang đo và mức độ tác động của từng yếu tố.
\item Khám phá khác biệt theo các đặc điểm nền được lưu trong tệp dữ liệu: giới tính, khu vực cư trú và thu nhập gia đình.
\item Đề xuất hàm ý quản trị cho trường đại học, trường THPT, gia đình và học sinh.
\end{enumerate}

\section{Câu hỏi nghiên cứu}
\begin{enumerate}[label=RQ\arabic*:]
\item Những yếu tố nào cấu thành quyết định chọn trường/ngành của học sinh THPT tại ĐBSCL?
\item Chiều hướng và mức độ tác động của mỗi yếu tố ra sao?
\item Mức độ quyết định có khác nhau giữa các nhóm học sinh hay không?
\item Các bên liên quan nên ưu tiên giải pháp nào?
\end{enumerate}

\section{Đối tượng và phạm vi nghiên cứu}
Đối tượng nghiên cứu là mối quan hệ giữa các yếu tố cá nhân, nghề nghiệp, cơ sở đào tạo, tài chính, xã hội, truyền thông và địa lý với quyết định chọn trường/ngành. Đối tượng khảo sát là học sinh THPT tại Đồng bằng sông Cửu Long trong năm học 2025--2026. Tệp phân tích cuối cùng gồm 486 quan sát. Việc suy rộng được giới hạn trong đặc điểm của mẫu và các biến được lưu trong tệp dữ liệu.

\section{Phương pháp nghiên cứu}
\subsection{Thiết kế tổng thể}
Nghiên cứu sử dụng thiết kế định lượng cắt ngang. Bảng hỏi được xây dựng từ cơ sở lý thuyết và các nghiên cứu liên quan, sau đó diễn đạt phù hợp với bối cảnh học sinh THPT. Các biến quan sát sử dụng thang Likert 5 mức từ 1 ``hoàn toàn không đồng ý'' đến 5 ``hoàn toàn đồng ý''.

\subsection{Chọn mẫu và cỡ mẫu}
Cỡ mẫu tối thiểu theo công thức Cochran ở mức tin cậy 95\%, sai số 5\%, $p=0{,}5$ là
\begin{equation}n_0=\frac{z^2p(1-p)}{e^2}=\frac{1{,}96^2\times0{,}5\times0{,}5}{0{,}05^2}\approx385.\end{equation}
Cỡ mẫu phân tích cuối cùng là 486, vượt ngưỡng tối thiểu theo công thức Cochran và đáp ứng yêu cầu thực hiện EFA cho 28 biến độc lập. Do tệp phân tích không lưu mã trường và tỉnh/thành, nghiên cứu không ước lượng trọng số mẫu và không khẳng định tính đại diện xác suất cho toàn vùng.

\subsection{Quy trình xử lý}
Tệp phân tích được kiểm tra số quan sát, kiểu dữ liệu, khoảng giá trị và giá trị thiếu trước khi tính điểm thang đo. Sau thống kê mô tả, nghiên cứu kiểm định Cronbach's Alpha, EFA, tương quan Pearson và hồi quy tuyến tính đa biến. Kiểm định Welch/ANOVA được dùng khi so sánh nhóm. Ngưỡng ý nghĩa là 5\%. Toàn bộ kết quả chính có thể được tính lại bằng mã trong thư mục \texttt{analysis/}.

\section{Đóng góp của nghiên cứu}
Về học thuật, nghiên cứu tích hợp lựa chọn trường và ngành trong cùng một mô hình, nhấn mạnh tính phù hợp cá nhân bên cạnh tín hiệu thị trường. Về thực tiễn, kết quả giúp ĐHCT và các trường trong vùng thiết kế nội dung tuyển sinh minh bạch, phân nhóm học sinh và phối hợp với trường THPT. Bộ câu hỏi ở phụ lục có thể tái sử dụng sau khi được thẩm định.

\section{Kết cấu luận văn}
Ngoài các phần đầu, kết luận, tài liệu tham khảo và phụ lục, nội dung chính gồm ba chương theo quy định được cung cấp: Tổng quan; Cơ sở lý thuyết; Nội dung và kết quả nghiên cứu.

\section{Đạo đức nghiên cứu và bảo vệ dữ liệu}
Phần giới thiệu của bảng hỏi nêu rõ tính tự nguyện, quyền dừng trả lời và nguyên tắc báo cáo tổng hợp. Tệp phân tích không chứa họ tên hoặc thông tin liên lạc. Với người chưa đủ 18 tuổi, quy trình đồng thuận cần tuân theo yêu cầu của cơ sở đào tạo và trường THPT. Dữ liệu được mã hóa bằng số thứ tự quan sát và chỉ sử dụng cho mục đích phân tích nghiên cứu.

\section{Khung tiến độ và nguồn lực}
\begin{table}[H]\centering\caption{Tiến độ thực hiện nghiên cứu}\begin{tabularx}{\textwidth}{p{3cm}Xp{3cm}}\toprule Giai đoạn & Công việc & Sản phẩm\\\midrule Tháng 1 & Lược khảo và hoàn thiện đề cương & Đề cương nghiên cứu\\Tháng 2 & Hiệu chỉnh thang đo và bảng hỏi & Bảng hỏi khảo sát\\Tháng 3--5 & Thu thập và kiểm tra dữ liệu & Tệp dữ liệu phân tích\\Tháng 6 & Phân tích và kiểm tra độ bền & Bảng và hình kết quả\\Tháng 7 & Viết, rà soát và hoàn thiện & Luận văn và phụ lục\\\bottomrule\end{tabularx}\end{table}
